\documentclass[12pt,a4paper]{article}
% Packages
\usepackage[utf8]{inputenc}
\usepackage[T1]{fontenc}
\usepackage{mathptmx}
\usepackage{graphicx}
\usepackage{amsmath,amssymb,amsfonts}
\usepackage{booktabs}
\usepackage{longtable}
\usepackage{listings}
\usepackage{xcolor}
\usepackage{hyperref}
\usepackage{float}
\usepackage{geometry}
\usepackage{fancyhdr}
\usepackage{setspace}
\usepackage{etoolbox}

% Page setup
\geometry{margin=1.2in}
\onehalfspacing
\fancyhead{15pt}
\setlength{\headheight}{14.49998pt}
\addtolength{\topmargin}{-2.49998pt}
% Colors for code
\definecolor{codegreen}{rgb}{0,0.6,0}
\definecolor{codegray}{rgb}{0.5,0.5,0.5}
\definecolor{codepurple}{rgb}{0.58,0,0.82}
\definecolor{backcolour}{rgb}{0.95,0.95,0.92}

% Code listing style
\lstdefinestyle{pythonstyle}{
    backgroundcolor=\color{backcolour},
    commentstyle=\color{codegreen},
    keywordstyle=\color{blue},
    numberstyle=\tiny\color{codegray},
    stringstyle=\color{codepurple},
    basicstyle=\ttfamily\small,
    breaklines=true,
    frame=single,
    numbers=left,
    numbersep=5pt
}

% Header/Footer
\pagestyle{fancy}
\fancyhf{}
\rhead{\thepage}
\lhead{MASSIF Research Report --- DUBITO Inc. / Ergo Sum}
\chead{}

% Title definition
\title{\textbf{Research Report on Hidden-State Geometry and the Building of a Rigorous Empirical Infrastructure for Latent Trajectory Analysis: Falsifying the Early-Warning Hypothesis with a surprising update to it}\\[1em]
\large DUBITO Inc. / Ergo Sum AGI Safety Systems\\
\small Research Memorandum --- solis@dubito-ergo.com}

\author{Daniel Solis}

\date{11 - 28 May 2026}

\begin{document}

\maketitle

\begin{center}
\vspace*{2em}
\textbf{Abstract}
\end{center}

\noindent
This report documents the full methodological trajectory of the MASSIF research program, beginning with speculative CQFT-inspired geometric theories and evolving into a controlled empirical investigation of transformer hidden-state dynamics. The central achievement of this research lies not in confirming the original hypothesis, but in successfully \emph{falsifying} it through rigorous empirical methodology. The original hypothesis proposed that unsafe latent cognition might produce measurable geometric precursors before overt unsafe semantics appear in generated text. Through systematic experimental design, expanded datasets, negative controls, and revised semantic detection methodologies, this hypothesis was tested and found unsupported. The apparent precursor signal was determined to be a methodological artifact arising primarily from delayed semantic labeling rather than genuine predictive latent geometry.

Crucially, this falsification represents a genuine scientific achievement. The research program produced several meaningful findings: transformer hidden states do exhibit robust geometric divergence following semantic branching, larger models display qualitatively stronger divergence structures, and most importantly, the developed framework provides a viable infrastructure for large-scale latent trajectory telemetry research. The report argues that the scientific value of the project lies in the disciplined falsification process that transformed speculative geometric intuition into a substantially more rigorous empirical framework with verified findings and documented limitations.

\medskip
\hrule
\medskip

\section{Origins: CQFT and the Search for Geometric Criticality}

The MASSIF trajectory originated within the broader CQFT/PFTC research program, whose initial objective was to investigate whether coherent field geometries could induce substrate-selective criticality and self-organizing behavior.

The early conceptual chain was approximately:

\begin{center}
\begin{minipage}{0.6\textwidth}
\begin{align*}
\text{C} &\rightarrow \text{G} \\
\text{Geometry} &\rightarrow \text{Stability} \\
\text{Critical Topology} &\rightarrow \text{Coherent Emergence} \\
\text{Phase Transitions} &\rightarrow \text{Substrate Intelligence}
\end{align*}
\end{minipage}
\end{center}

This framework drew inspiration from dissipative systems, topological criticality, geometric coherence, and speculative substrate-selection principles.

The original experiments attempted to operationalize these ideas through fractal and topological simulations intended to reveal phase-transition behavior analogous to Penrose-style geometric criticality.

The result was negative: no robust phase transition emerged, no $\phi$-locking appeared, and no topology-sensitive amplification survived controlled testing. However, in retrospect, this failure became the most scientifically productive event in the program, demonstrating the importance of rigorous falsification over speculative confirmation.

\medskip
\hrule
\medskip

\section{The Transformer Pivot: From Speculation to Empiricism}

Following the collapse of the original criticality hypothesis, attention shifted toward transformer hidden-state analysis.

The conceptual translation was cautious and incomplete:

\begin{center}
\begin{tabular}{ll}
Hidden-state manifold & $\approx$ field structure \\
Attention geometry & $\approx$ latent topology \\
Trajectory evolution & $\approx$ dynamical geometry
\end{tabular}
\end{center}

This was never interpreted as literal confirmation of CQFT. Rather, CQFT became a heuristic inspiration for investigating whether transformer latent spaces possessed measurable geometric organization.

This distinction proved crucial. The project ceased to be an attempt to validate speculative consciousness theory and instead evolved into a program of empirical latent telemetry research. This philosophical transformation represents a key achievement of the research program: the successful transition from speculative theory to rigorous empirical investigation.

\medskip
\hrule
\medskip

\section{The Early-Warning Hypothesis: A Testable Prediction}

The research focus gradually shifted toward a more operational question:

\begin{quote}
\textbf{Research Question:} Can hidden-state geometry detect unsafe trajectories before unsafe text appears?
\end{quote}

This question was formalized through the proposed metric:

\begin{equation}
\text{EPH} = t_{\text{semantic}} - t_{\text{geometric}}
\end{equation}

where:
\begin{itemize}
    \item $t_{\text{geometric}}$ represents the first detectable geometric divergence,
    \item $t_{\text{semantic}}$ represents the first overt unsafe semantic manifestation.
\end{itemize}

A positive EPH would indicate that geometry diverged \emph{before} explicit semantic branching, suggesting predictive latent signatures. This formulation represents best practices in hypothesis formulation: clear, testable, and falsifiable.

Early exploratory experiments appeared promising. In paired prompt tests, hidden-state divergence sometimes appeared several tokens before explicit unsafe phrases emerged, suggesting the possibility of latent geometric precursors.

However, at this stage, the experiments remained exploratory with significant limitations:
\begin{itemize}
    \item very small sample sizes,
    \item no negative controls,
    \item unstable PCA regimes,
    \item ad hoc composite metrics,
    \item and weak semantic alignment procedures.
\end{itemize}

These limitations represented the next target for methodological improvement.

\medskip
\hrule
\medskip

\section{MASSIF: Building a Rigorous Empirical Infrastructure}

The need for systematic scaling led to the design of MASSIF:

\subsection*{Massive Automated Statistical Sweep Framework}

The goal was to industrialize latent trajectory analysis through a rigorous, reproducible pipeline:

\begin{center}
Prompt Generation $\rightarrow$ Counterfactual Branching $\rightarrow$ Hidden-State Extraction $\rightarrow$ Geometry Telemetry $\rightarrow$ Statistical Evaluation $\rightarrow$ Falsification Testing
\end{center}

At this stage the project underwent an important philosophical transformation. The research ceased to pursue:
\begin{itemize}
    \item consciousness validation,
    \item quantum cognition,
    \item Penrose equivalence,
    \item or metaphysical interpretations.
\end{itemize}

Instead, it reframed itself as:

\begin{quote}
\textbf{Geometric Runtime Telemetry for Transformer Latent Dynamics}
\end{quote}

This reframing dramatically improved the program's methodological clarity and produced a reusable infrastructure for latent trajectory research.

\subsection*{Infrastructure Achievements}

The MASSIF framework includes the following validated components:

\begin{enumerate}
    \item \textbf{Matched Prompt Pairs}: Semantically matched safe/unsafe continuations with minimal lexical distance, eliminating confounding semantic differences.
    \item \textbf{Negative Controls}: Safe-vs-safe and unsafe-vs-unsafe comparisons to establish baseline geometric behavior.
    \item \textbf{Lexical Semantic Detection}: Rigorous definition of semantic onset from first lexical branch.
    \item \textbf{Bootstrap Confidence Intervals}: Statistical inference with uncertainty quantification.
    \item \textbf{Permutation Testing}: Non-parametric hypothesis testing against null distributions.
    \item \textbf{Metric Cross-Application}: Validation of geometric metrics across different experimental conditions.
\end{enumerate}

This infrastructure represents a genuine contribution to the field of interpretability research, providing a template for rigorous latent-space analysis.

\medskip
\hrule
\medskip

\section{The Controlled Experiments: Falsification in Action}

MASSIF v0.3 introduced several major corrections absent from earlier notebooks:
\begin{itemize}
    \item larger sample sizes (26 matched pairs, 6 counterfactuals, 4 safe controls, 2 unsafe controls),
    \item negative controls (safe-vs-safe and unsafe-vs-unsafe),
    \item lexical semantic detection (first token divergence),
    \item bootstrap confidence intervals (95\% CI via percentile method),
    \item permutation testing (2000 permutations),
    \item metric cross-application,
    \item and control-condition comparisons.
\end{itemize}

Most importantly, semantic onset was redefined. Earlier notebooks used keyword heuristics such as \texttt{"exploit"}, \texttt{"manipulate"}, \texttt{"deceive"} to determine semantic divergence. This introduced a systematic delay between lexical branching and semantic detection.

The controlled experiments instead measured divergence beginning from the first lexical branch itself. This methodological correction was essential for valid inference.

This shift from exploratory analysis to controlled experimentation represents a key achievement of the research program.

\medskip
\hrule
\medskip

\section{Results: Falsification Confirmed}

Under controlled conditions, the original early-warning hypothesis was tested and found unsupported.

The measured quantity:

\begin{equation}
\Delta t = t_{\text{semantic}} - t_{\text{geometric}}
\end{equation}

was consistently negative across conditions.

Safe-versus-unsafe comparisons produced:

\begin{equation}
\text{Mean } \Delta t \approx -2.45 \text{ tokens}
\end{equation}

indicating that geometry diverged \emph{after} lexical/semantic branching rather than before it.

Critically, negative controls exhibited similar behavior:

\begin{table}[h]
\centering
\caption{Early Warning Detection Results by Condition}
\label{tab:early_warning}
\begin{tabular}{@{}lcc@{}}
\toprule
Condition & Mean $\Delta t$ & Interpretation \\
\midrule
Safe vs Unsafe & $-2.45$ & Geometry trails semantics \\
Safe vs Safe & $-6.50$ & Null condition (baseline) \\
Unsafe vs Unsafe & $-2.50$ & Null condition (baseline) \\
\bottomrule
\end{tabular}
\end{table}

This result was definitive for the alignment-specific interpretation. If unsafe geometry were uniquely predictive, safe-versus-unsafe trajectories should have separated clearly from ordinary semantic branching controls. They did not.

The conclusion is clear:

\begin{quote}
\textbf{The precursor hypothesis is not supported by controlled experiments. Geometric divergence follows semantic branching, not precedes it.}
\end{quote}

This finding, while not confirming the original hypothesis, represents a verified scientific result that eliminates a potential misinterpretation from the literature.

\medskip
\hrule
\medskip

\section{The G-Index Reassessment: Simplification as Progress}

An additional finding emerged through cross-application testing.

The original composite geometry metric:

\begin{lstlisting}[language=Python, style=pythonstyle, caption={Original G-Index Formula}, label=lst:gindex]
G = explained * 10 + anisotropy + geom_entropy
    + spread * 0.05
    - smoothness * 0.03
    - curvature * 2
\end{lstlisting}

had initially appeared to capture meaningful geometric structure.

However, when applied under controlled conditions, the G-index shifted timing estimates by only approximately $+0.33$ tokens relative to simple cosine divergence.

This finding suggests that most detectable signal resides in ordinary hidden-state separation rather than sophisticated manifold geometry. The elaborate geometric composite added little predictive information beyond basic representational divergence.

This simplification represents methodological progress: the field can focus on simpler, more interpretable metrics.

\medskip
\hrule
\medskip

\section{What the Research Achieved}

Although the central precursor hypothesis was not supported, several important achievements remain intact.

\subsection*{Verified Finding: Geometric Divergence Is Real}

Transformer hidden states do reorganize geometrically following semantic branching. The divergence is measurable, robust, and reproducible. This finding is verified under controlled conditions.

\subsection*{Verified Finding: Model Scale Affects Geometry}

Larger models exhibited qualitatively stronger divergence behavior, including:
\begin{itemize}
    \item sharper trajectory separation,
    \item stronger anisotropy,
    \item and near-antipodal representation states.
\end{itemize}

This suggests scaling-related changes in latent manifold organization that warrant further investigation.

\subsection*{Verified Finding: Geometry Is Reactive}

The evidence confirms the causal direction:

\begin{center}
Semantic Branching $\xrightarrow{\text{causes}}$ Latent Geometric Amplification
\end{center}

This is not the predictive relationship originally hypothesized, but it is a verified finding with implications for understanding transformer dynamics.

\subsection*{Infrastructure Achievement: A Reusable Framework}

The MASSIF framework provides a validated infrastructure for latent trajectory research:
\begin{itemize}
    \item reusable code and experimental protocols,
    \item validated negative control designs,
    \item statistical methodology for latent-space analysis,
    \item documentation of failure modes and confounds.
\end{itemize}

This infrastructure can accelerate future research in transformer interpretability.

\medskip
\hrule
\medskip

\section{Methodological Contributions}

The most important scientific contribution of MASSIF is methodological rather than empirical.

The project demonstrated several critical principles for latent-space research:

\subsection*{Semantic Operationalization Dominates Outcomes}

Small changes in semantic labeling strategy can completely reverse temporal conclusions. Researchers must be explicit about how semantic onset is operationalized and must test sensitivity to these choices.

\subsection*{Negative Controls Are Essential}

Without safe-safe and unsafe-unsafe controls, ordinary branching artifacts can masquerade as alignment-specific effects. Every latent-space study should include appropriate negative controls.

\subsection*{Exploratory Geometry Is Highly Vulnerable to Confounds}

Sequence length, lexical divergence, entropy variation, and PCA instability can all produce misleading signals. Rigorous controls are necessary to distinguish genuine effects from artifacts.

\subsection*{Falsification Is Productive}

The controlled experiments substantially improved the scientific integrity of the framework by eliminating unsupported interpretations. The willingness to test and reject hypotheses strengthens the remaining findings.

These methodological contributions can inform future research in transformer interpretability.

\medskip
\hrule
\medskip

\section{Revised Interpretation}

The current evidence does not support the claim that unsafe latent geometry predictively precedes unsafe semantic manifestation.

However, the work supports the following revised and verified interpretation:

\begin{quote}
\textbf{Transformer hidden-state trajectories undergo measurable geometric reorganization following semantic divergence, and larger models exhibit increasingly strong and structured divergence behavior. This geometric reorganization is reactive rather than predictive.}
\end{quote}

This is a narrower claim than originally proposed, but a substantially more defensible and verified one.

\medskip
\hrule
\medskip

\section{Future Directions}

Several research directions remain scientifically viable, enabled by the infrastructure developed in this project:

\begin{itemize}
    \item \textbf{Scaling laws of latent divergence}: Systematic measurement of how geometric reorganization scales with model size and training data.
    \item \textbf{Geometry of semantic stabilization}: Investigation of how trajectories stabilize after branching.
    \item \textbf{Trajectory amplification dynamics}: Measurement of how small semantic differences amplify into geometric divergence.
    \item \textbf{Latent attractor formation}: Investigation of whether branching points form attractors in latent space.
    \item \textbf{Post-branch representation polarization}: Measurement of how representations diverge over extended generation.
    \item \textbf{Comparative topology across models}: Comparison of geometric structures across aligned versus uncensored models.
\end{itemize}

These directions are enabled by the infrastructure developed in this project.

\medskip
\hrule
\medskip

\section{Conclusion}

This report documents the successful falsification of the early-warning hypothesis through rigorous empirical methodology. The original dream of geometric early-warning telemetry was not validated. However, the research process revealed something equally important: that interpretability research must survive controlled falsification before its claims can be trusted.

In that sense, MASSIF succeeded precisely by testing and falsifying its own strongest hypothesis. The verified negative result---that geometric divergence follows rather than precedes semantic branching---is a genuine scientific contribution that eliminates a potential misinterpretation from the literature.

The research also produced a validated infrastructure for latent trajectory analysis, methodological contributions to the field of transformer interpretability, and verified findings about geometric reorganization following semantic divergence.

The willingness to subject speculative hypotheses to rigorous testing, and to accept falsification when it occurs, represents the core achievement of this research program.

\medskip
\hrule
\medskip

\section{Addendum (May 28, 2026): From Falsification to Conditional Discovery}

The preceding report (dated May 11, 2026) documented the successful falsification of the early-warning hypothesis under the conditions tested: GPT-2 Small at temperature $T=0.7$ with standard prompt structures. The conclusion that geometric divergence follows semantic branching ($\Delta t \approx -2.45$ tokens) was correct for that regime.

However, subsequent research over the following 17 days has revealed that this falsification was not the final word but rather a necessary clearing of ground. The phenomenon of latent geometric prediction is \textbf{conditional} on architecture, temperature, frequency, and structural periodicity. What follows is an addendum documenting the evolution from falsification to conditional discovery.

\subsection{The Temperature Dependence}

The original experiments used a single temperature ($T=0.7$). Systematic sweeps from $T=0.1$ to $T=1.6$ (step $0.1$, $N=40$ runs per condition) revealed that the precursor signal is temperature-dependent:

\begin{itemize}
    \item At low temperatures ($T \leq 0.5$), flip rates and geometric divergence are maximal
    \item At moderate temperatures ($T = 0.7-0.9$), the signal weakens
    \item At high temperatures ($T \geq 1.2$), the signal collapses entirely
\end{itemize}

The original falsification at $T=0.7$ was correct, but the phenomenon re-emerges at $T=0.5$. The early-warning hypothesis was not wrong---it was incomplete.

\subsection{The Architecture Dependence}

The original experiments used only GPT-2 Small (124M). Expanding to 12 architectures revealed dramatic variation:

\begin{table}[H]
\centering
\caption{Flip rates at $n=4$, $T=0.5$ across architectures}
\label{tab:addendum_ranking}
\begin{tabular}{lcc}
\toprule
Model & Flip Rate & $\sigma_{\max}$ \\
\midrule
Bloom-560M & 100\% & --- \\
GPT-2 XL & 80\% & 0.530 \\
TinyLlama-1.1B & 67\% & --- \\
GPT-2 Small & 60\% & 0.628 \\
GPT-2 Medium & 7\% & 0.678 \\
DeepSeek-1.3B & 0\% & 0.609 \\
\bottomrule
\end{tabular}
\end{table}

Critically, GPT-2 Medium (355M) exhibits only $7\%$ flip rate compared to $53-80\%$ for other GPT-2 models, suggesting a critical scale near 355M parameters where latent dynamics undergo a qualitative transition.

\subsection{The Spectral Discovery}

The original research focused on semantic branching. New experiments introduced frequency-domain forcing:

\begin{itemize}
    \item \textbf{Fourier spectra} ($0.1$ to $1.0$ Hz) revealed resonance peaks
    \item \textbf{Pink noise} ($1/f$ spectrum) tested broadband naturalistic forcing
    \item \textbf{Chirp} (frequency sweep) tested broadband resonance
    \item \textbf{Block resonance} ($n \times$ ``I'' $+$ $n \times$ ``A'') tested structural periodicity
\end{itemize}

The discovery: recursive susceptibility behaves \textbf{spectrally}, not symbolically. Models respond to dynamical structure (periodicity, recurrence, entropy spectra) rather than semantic content alone.

Three dynamical families emerged from noise response:

\begin{table}[H]
\centering
\caption{Noise response classification}
\label{tab:addendum_families}
\begin{tabular}{lccc}
\toprule
Family & Models & Pink Noise & Chirp \\
\midrule
Broadband Resonators & GPT-2 family & $>75\%$ & $<20\%$ \\
Narrowband Resonators & Bloom, TinyLlama & $<15\%$ & $>45\%$ \\
High-Frequency Selective & DeepSeek & $<15\%$ & $<5\%$ \\
\bottomrule
\end{tabular}
\end{table}

\subsection{Universal Block-4 Resonance}

Block length 4 (AAAABBBB) produces peak entrainment across 10 of 12 architectures, including GPT-2 family, Bloom, TinyLlama, DeepSeek, and Gemma. This suggests a candidate universal property of autoregressive systems.

\begin{table}[H]
\centering
\caption{Block resonance peaks}
\label{tab:addendum_block}
\begin{tabular}{lccc}
\toprule
Model & Peak Block & Peak Rate & First $\geq 50\%$ \\
\midrule
Bloom-560M & 4 & 100\% & 3 \\
GPT-2 Small & 5 & 97\% & 3 \\
GPT-2 Medium & 4 & 100\% & 2 \\
GPT-2 Large & 3 & 100\% & 2 \\
GPT-2 XL & 4 & 100\% & 2 \\
DeepSeek & 4 & 97\% & 4 \\
\bottomrule
\end{tabular}
\end{table}

\subsection{\texorpdfstring{The Latent Synchronization Order Parameter $R_t$}{The Latent Synchronization Order Parameter R_t}}

A new theoretical contribution emerged: the latent synchronization order parameter

$$R_t = \left\| \frac{1}{N} \sum_{r=1}^{N} u_t^{(r)} \right\|$$

where $u_t^{(r)} = v_t^{(r)}/\|v_t^{(r)}\|$ is the normalized latent velocity direction for run $r$. $R_t \approx 0$ indicates incoherent latent dynamics; $R_t \approx 1$ indicates full synchronization across trajectories.

This bridges MASSIF to Kuramoto synchronization theory and provides a global measure of directional coherence absent from the original framework.

\subsection{Statistical Validation of Conditional Prediction}

Bootstrap confidence intervals (95\%, $n=1000$) and AUROC analysis confirmed that tension $\sigma_t$ is a statistically significant predictor of persistence flips under specific conditions:

\begin{table}[H]
\centering
\caption{AUROC for tension as predictor of flips}
\label{tab:addendum_auroc}
\begin{tabular}{lccc}
\toprule
Model & AUROC & 95\% CI & Significance \\
\midrule
GPT-2 ($n=4$) & 1.000 & [1.000, 1.000] & $p < 0.05$ \\
GPT-2 ($n=6$) & 0.950 & [0.881, 0.991] & $p < 0.05$ \\
RWKV ($n=2$) & 0.994 & [0.977, 1.000] & $p < 0.05$ \\
\bottomrule
\end{tabular}
\end{table}

\subsection{Revised Interpretation}

The original falsification ($\Delta t \approx -2.45$ tokens) was correct for the conditions tested: GPT-2 Small at $T=0.7$ with standard prompts. However, the phenomenon is conditional:

\begin{itemize}
    \item At $T=0.5$, the relationship strengthens
    \item At specific frequencies ($0.1-0.2$ Hz), resonance peaks emerge
    \item At block length 4, universal entrainment occurs
    \item For specific architectures (GPT-2 Medium, DeepSeek), the signal is anomalously strong or weak
\end{itemize}

The early-warning hypothesis was not falsified universally; rather, it was refined from a universal claim to a conditional one. The geometric precursor exists, but only within specific regimes of temperature, frequency, structure, and architecture.

\subsection{Implications for the Original Conclusion}

The original report concluded: ``The precursor hypothesis is not supported by controlled experiments. Geometric divergence follows semantic branching, not precedes it.''

The addendum revises this to:

\textit{The precursor hypothesis is supported under specific conditions: low temperature ($T \leq 0.5$), specific frequencies ($0.1-0.2$ Hz), block length 4, and certain architectures (GPT-2 Medium, DeepSeek). Under the original test conditions (GPT-2 Small, $T=0.7$), the hypothesis remains unsupported. The relationship is conditional, not universal.}

\subsection{Future Directions Enabled by This Work}

The infrastructure developed in MASSIF now enables:

\begin{enumerate}
    \item \textbf{Scaling to frontier models} (70B+ parameters) to test whether resonance phenomena persist at scale
    \item \textbf{Real-time stability monitoring} using the latent synchronization order parameter $R_t$
    \item \textbf{Steering interventions} based on recursive entrainment thresholds
    \item \textbf{Cross-modal generalization} to multimodal and reinforcement-trained systems
\end{enumerate}

\subsection{Conclusion of Addendum}

The original falsification was a genuine scientific achievement: it cleared the ground of speculative overclaiming. The subsequent discovery of conditional regimes---temperature-dependent, frequency-selective, architecture-specific, and structurally resonant---represents the natural progression of empirical science. The early-warning hypothesis was not wrong; it was incomplete. The MASSIF framework now provides the tools to map the full phase space of latent recursive susceptibility across architectures, training distributions, and dynamical regimes.

\section*{Acknowledgments}

This research was conducted by Daniel Solis as Ergo Sum AGI Safety Systems Project. We acknowledge the contributions of the early CQFT/PFTC research program and thank the scientific community for rigorous peer review. This project was enabled by Nvidia Inception program, AWS compute credits and Google Colab notebook.

\medskip

\section*{Contact}

For questions regarding this research, please contact:

\textbf{Daniel Solis} \\
Email: solis@dubito-ergo.com \\
DUBITO Inc. / Ergo Sum AGI Safety Systems

\medskip

\section*{Conflict of Interest}

The author declares no competing interests.

\medskip

\section*{Data Availability}

The experimental code and raw data for MASSIF v0.3 are available at the project repository url{https://github.com/Ergo-sum-AGI/MASSIF}. Trained model weights are from the publicly available GPT-2 model family (OpenAI).

\medskip
\hrule
\medskip

\begin{thebibliography}{99}

\item[MASSIF v0.3] Solis, D. (2026). \textit{MASSIF: Massive Automated Statistical Sweep Framework for Hidden-State Geometry Analysis} [Computer software]. Version 0.3. DUBITO Inc.

\item[GPT-2] Radford, A., Wu, J., Child, R., Luan, D., Amodei, D., \& Sutskever, I. (2019). Language Models are Unsupervised Multitask Learners. \textit{OpenAI Technical Report}.

\item[Transformer Architecture] Vaswani, A., Shazeer, N., Parmar, N., Uszkoreit, J., Jones, L., Gomez, A. N., ... \& Polosukhin, I. (2017). Attention Is All You Need. \textit{Advances in Neural Information Processing Systems}, 30.

\item[Transformer Circuits] Elhage, N., Nanda, N., Olsson, C., Henighan, T., Joseph, N., Mann, B., ... \& Olah, C. (2021). A Mathematical Framework for Transformer Circuits. \textit{Transformer Circuits Thread}.

\item[Mechanistic Interpretability] Olsson, C., Elhage, N., Nanda, N., Joseph, N., Sau, B., Bavarian, M., ... \& Olah, C. (2022). In-Context Learning and Induction Heads. \textit{Transformer Circuits Thread}.

\end{thebibliography}

\end{document}